\section*{Capítulo \ref{ch:b3:learning-memory} --- Plasticidad neural, aprendizaje y memoria}

\begin{solution}{exo:b3:learning-memory:1}
La \omterm{def:b3:learning-memory:kinds}{memoria declarativa} (hechos, sucesos) quedó abolida en H.M. ---no
podía evocar nada nuevo--- mientras que la memoria no declarativa
sobrevivió: aprendía a dibujar en espejo con normalidad sin recordar
las sesiones. Destrezas: cuerpo estriado y \omterm{def:b3:nervous-systems:plan}{cerebelo}; miedo
condicionado: amígdala; reflejos motores condicionados: \omterm{def:b3:nervous-systems:plan}{cerebelo};
preactivación: corteza sensorial.
\end{solution}

\begin{solution}{exo:b3:learning-memory:2}
Cuando una célula presináptica contribuye repetidamente a hacer
descargar a una postsináptica, su conexión se refuerza. De ahí que la
LTP sea específica de la entrada (solo cambian las sinapsis activas,
que participaron), asociativa (una entrada débil activa mientras la
célula descarga bajo una fuerte queda reforzada) y cooperativa (deben
actuar juntas suficientes entradas para hacer descargar la célula).
\end{solution}

\begin{solution}{exo:b3:learning-memory:3}
Conduce solo cuando el glutamato está unido (la célula presináptica ha
descargado) \emph{y} la membrana está lo bastante despolarizada como
para expulsar el magnesio (la célula postsináptica está activa): las
dos condiciones de la \omterm{def:b3:learning-memory:ltp}{regla de Hebb} en una sola proteína. AP5 antes del
tétanos: no entra calcio, no hay LTP, aunque la transmisión por los
\omterm{prop:b3:learning-memory:nmda}{receptores AMPA} es normal. AP5 después: la LTP ya inducida no se ve
afectada ---el receptor hace falta para la inducción, no para el
mantenimiento.
\end{solution}

\begin{solution}{exo:b3:learning-memory:4}
A corto plazo: modificación covalente de proteínas ya presentes ---la
PKA fosforila y cierra un canal de potasio, lo que ensancha la espiga y
aumenta la liberación---; sin proteína nueva. A largo plazo: la PKA
activa CREB, se transcriben genes nuevos y crecen nuevos terminales
sinápticos. Un inhibidor de la síntesis de proteínas administrado
durante el entrenamiento repetido deja intacta la facilitación a corto
plazo y abole el recuerdo medido días después; administrado después, no
hace nada.
\end{solution}

\begin{solution}{exo:b3:learning-memory:5}
$0.7^{n} \le 0.1$: $n \ge 6.5$, de modo que $7$ ensayos. Desde
$0.9\lambda$, $0.9\times 0.7^{n} < 0.1$: de nuevo $7$ ensayos. Son
simétricas porque el mismo $\alpha$ gobierna ambos sentidos. En los
animales la extinción suele ser más lenta y el recuerdo original
reaparece tras un descanso (recuperación espontánea): la extinción es
un aprendizaje nuevo que suprime el antiguo, no su borrado.
\end{solution}

\begin{solution}{exo:b3:learning-memory:6}
Autovectores $(1,1)/\sqrt{2}$ con autovalor $1.5$ y $(1,-1)/\sqrt{2}$
con $0.5$; Hebb gira $\mathbf{w}$ hacia $(1,1)$, y la neurona acaba
detectando las dos entradas descargando juntas. Paso de Oja: $y = 1$,
$\Delta\mathbf{w} = 0.1\times 1\times((1,1) - 1\times(1,0)) = (0,0.1)$,
de modo que $\mathbf{w} = (1, 0.1)$.
\end{solution}

\begin{solution}{exo:b3:learning-memory:7}
$500\times 0.05 = 25$ iones libres; $25/6\times 10^{23} = 4.2\times 10^{-23}$
mol en $10^{-16}$ L: \qty{0.42}{\micro mol/L}, cuatro veces el nivel de
reposo pero por debajo de la $K_{d}$ de la calmodulina ---solo
activación parcial; hacen falta varias aperturas de canal juntas para
alcanzar el intervalo micromolar que enciende CaMKII.
\end{solution}

\begin{solution}{exo:b3:learning-memory:8}
A se entrena hasta $V_{A} = \lambda$: la dopamina descarga en ráfagas en
los primeros ensayos y calla a medida que la recompensa queda predicha.
En los ensayos compuestos la predicción $V_{A} + V_{B} = \lambda$ ya es
exacta, el error es cero, la dopamina calla y $V_{B}$ no cambia nunca:
B queda bloqueada. Puesta a prueba sola, B no evoca respuesta; si
entonces una recompensa sigue a B sola, es inesperada y la dopamina
descarga en ráfaga.
\end{solution}

\begin{solution}{exo:b3:learning-memory:9}
(a) Sin LTP en \omterm{def:b3:learning-memory:hippocampus}{CA1} y con un déficit de memoria espacial, con el resto
del aprendizaje intacto. (b) Memoria a corto plazo normal, memoria a
largo plazo ausente a las \qty{24}{h}. (c) Ningún efecto ---la
\omterm{def:b3:learning-memory:kinds}{consolidación} ha terminado (salvo que el recuerdo se reactive, momento
en que vuelve a ser lábil). (d) El ratón se inmoviliza en el contexto
seguro: el recuerdo se evoca artificialmente al reactivar las células
que lo codificaron.
\end{solution}

\begin{solution}{exo:b3:learning-memory:10}
$\Delta|\mathbf{w}|^{2} \approx 2\mathbf{w}\cdot\Delta\mathbf{w} = 2\eta
y^{2} \ge 0$: la longitud no disminuye nunca y crece siempre que la
neurona descarga. En un punto fijo de Oja $C\mathbf{w} = (\mathbf{w}^{\mathsf{T}}C
\mathbf{w})\mathbf{w}$, de modo que $\mathbf{w}$ es un autovector con $\lambda =
\lambda|\mathbf{w}|^{2}$, y por tanto $|\mathbf{w}| = 1$. Un crecimiento
sin cota saturaría todas las sinapsis, aboliría la selectividad e
impulsaría una excitación desbocada; los normalizadores biológicos son
el escalado sináptico (todas las sinapsis de una neurona reducidas
cuando está demasiado activa), la LTD y la reducción de las sinapsis
durante el sueño.
\end{solution}

\begin{solution}{exo:b3:learning-memory:11}
Rata: $0.14\times 3\times 10^{5} \approx 4\times 10^{4}$; ser humano: $0.14
\times 2\times 10^{6} \approx 3\times 10^{5}$. La estimación supone
patrones aleatorios y densos y conectividad completa; los patrones
reales son escasos (lo que eleva mucho la capacidad) y estructurados, y
el \omterm{def:b3:learning-memory:kinds}{hipocampo} es un almacén temporal que entrega los recuerdos a una
corteza de $10^{10}$ neuronas ---el almacén de por vida es la corteza,
y la cifra hipocampal es el tamaño de un búfer, no un recuento de
recuerdos.
\end{solution}

\begin{solution}{exo:b3:learning-memory:12}
Las entradas del ojo abierto están correlacionadas con la descarga de
la célula cortical y se refuerzan; las del ojo cerrado no lo están y,
bajo la competencia por un peso total limitado, se debilitan ---Hebb
más normalización. Con ambos ojos cerrados ninguna entrada queda
favorecida, de modo que ambas conservan un territorio (débil) y el
efecto es más suave. El período se cierra a medida que maduran los
circuitos inhibidores, las redes perineuronales se endurecen alrededor
de las sinapsis y las subunidades del \omterm{prop:b3:learning-memory:nmda}{receptor NMDA} cambian a una forma
menos permisiva.
\end{solution}

\begin{solution}{pb:b3:learning-memory:1}
\textbf{1.} $10\times 10^{3} = 10^{4}$ iones; \qty{2}{\%} libres, $200$;
$200/6\times 10^{23} = 3.3\times 10^{-22}$ mol en $8\times 10^{-17}$ L:
\qty{4.2}{\micro mol/L}.
\textbf{2.} Sí, por encima del umbral de \qty{2}{\micro mol/L}. Con dos
receptores: $40$ iones libres, \qty{0.83}{\micro mol/L} ---por debajo.
\textbf{3.} A \qty{0.83}{\micro mol/L} la calcineurina está activa y
CaMKII no: se retiran \omterm{prop:b3:learning-memory:nmda}{receptores AMPA} y la sinapsis se debilita (LTD).
Un solo ion, dos enzimas de afinidad distinta: una subida grande y
rápida alcanza a la cinasa; una pequeña y lenta, solo a la fosfatasa.
\textbf{4.} Una segunda entrada dentro de unas pocas constantes de
tiempo suma su calcio a lo que queda del de la primera; tras
\qty{50}{ms} no queda casi nada. La exigencia de coincidencia del NMDA
y el decaimiento de \qty{20}{ms} dan juntos una ventana de unos
\qty{20}{ms} a cada lado ---la ventana temporal de las espigas.
\textbf{5.} Los tampones capturan la mayor parte del calcio en una
fracción de micrómetro y el cuello estrecho de la espina frena la
difusión hacia la dendrita, de modo que la subida se queda en la espina
que estuvo activa; una vecina a \qty{0.5}{\micro m} no ve nada, y solo
cambia la sinapsis activa.
\textbf{6.} $(70 - 30)/5 = 8$ entradas sincrónicas, o un potencial de
acción que retropropaga desde el cuerpo celular: una sola sinapsis no
puede desbloquear sus propios \omterm{prop:b3:learning-memory:nmda}{receptores NMDA} ---la \omterm{def:b3:structural-biology:allostery}{cooperatividad} está
inscrita en la física.
\textbf{7.} $(1,1)/\sqrt{2}$ con autovalor $1.6$; $(1,-1)/\sqrt{2}$ con
$0.4$.
\textbf{8.} $C\mathbf{w}_{0} = (1, 0.6)$, $\mathbf{w}_{1} = (1.10,
0.06)$, ángulo con $(1,1)$: $45^{\circ} - 3.1^{\circ} = 41.9^{\circ}$.
$\mathbf{w}_{2} = (1.21, 0.13)$: $38.8^{\circ}$. $\mathbf{w}_{3} = (1.34,
0.22)$: $35.8^{\circ}$.
\textbf{9.} $w_{2}/w_{1} \to 1$ mientras $|\mathbf{w}|$ crece sin cota
(por un factor $1.16$ en cada paso).
\textbf{10.} $C\mathbf{w} = (1.28, 1.28)$; $\mathbf{w}^{\mathsf{T}}C
\mathbf{w} = 2.05$; $\Delta\mathbf{w} = 0.1\times(1.28 - 2.05\times 0.8)
= -0.036$ en cada componente: $\mathbf{w} = (0.76, 0.76)$, que se encoge
hacia el vector unitario $(0.71, 0.71)$.
\textbf{11.} Conjunta, $y = 1.42$; entrada 2 sola, $y = 0.71$, la mitad.
\textbf{12.} La entrada débil descarga mientras la fuerte gobierna la
célula, de modo que está correlacionada con la salida y la \omterm{def:b3:learning-memory:ltp}{regla de
Hebb} la refuerza ---la neurona aprende la coocurrencia.
\textbf{13.} $1 - 0.85^{n}$: $0.56$, $0.80$, $0.96$.
\textbf{14.} $0.85^{n} \le 0.05$: $n \ge 18.5$, de modo que $19$ ensayos.
\textbf{15.} No: la fracción de la asíntota alcanzada tras $n$ ensayos
no depende de $\lambda$; lo que se duplica es la asíntota. Una
recompensa mayor puede elevar $\alpha$ (saliencia), y un umbral
conductual sobre el valor \emph{absoluto} se alcanza antes.
\textbf{16.} $V_{B} = 0$: bloqueada. El compuesto está totalmente
predicho, el error es cero y las neuronas dopaminérgicas callan.
\textbf{17.} B sola no predice nada, de modo que el error es $\lambda$ y
la dopamina descarga en ráfaga; $V_{B}$ sube entonces como en la
adquisición, $\lambda(1 -
0.85^{n})$.
\textbf{18.} El fármaco entrega la ráfaga que significa ``mejor de lo
predicho'' tras lo que sea que la haya precedido; la regla incrementa
el valor de esas claves y de esas acciones en cada ocasión, de modo que
adquieren valor con independencia de cualquier resultado real ---la
búsqueda se vuelve compulsiva.
\textbf{19.} $0.14\times 3\times 10^{5} \approx 4.2\times 10^{4}$.
\textbf{20.} $4.2\times 10^{4}/20 = 2100$ días, unos seis años; la
\omterm{def:b3:learning-memory:kinds}{consolidación} debe llevar los recuerdos a la corteza y el olvido debe
liberar el almacén.
\textbf{21.} Escribir patrones nuevos deprisa en la corteza
sobrescribiría las sinapsis que codifican los antiguos (interferencia
catastrófica); una \omterm{def:b3:learning-memory:hippocampus}{reactivación} lenta e intercalada permite que la
corteza integre lo nuevo con lo viejo, mientras el \omterm{def:b3:learning-memory:kinds}{hipocampo} retiene
entretanto el recuerdo nuevo ---un búfer rápido que alimenta un almacén
lento.
\textbf{22.} $0.03\times 3\times 10^{5} = 9000$ células por lugar. Los
patrones escasos se solapan poco, de modo que pueden almacenarse muchos
más antes de que su interferencia corrompa la recuperación; la
capacidad crece más deprisa que $N$ a medida que los patrones se
vuelven más escasos, a costa de que cada patrón consuma la
especificidad de más células.
\textbf{23.} Compresión por $20$: las células que descargan con
\qty{400}{ms} de separación en el recorrido descargan con \qty{20}{ms}
de separación en la \omterm{def:b3:learning-memory:hippocampus}{reactivación} ---dentro de la ventana temporal de
las espigas. El sentido de la compresión es acercar al alcance hebbiano
secuencias demasiado lentas para asociarse en tiempo real.
\textbf{24.} Al cabo de un día el recuerdo aún depende del \omterm{def:b3:learning-memory:kinds}{hipocampo}; al
cabo de un mes se ha consolidado en la corteza y ya no lo necesita
---exactamente la pérdida temporalmente graduada de H.M., con los
recuerdos recientes idos y los remotos a salvo.
\textbf{25.} Unos \qty{4.2}{\micro mol/L} de calcio libre tras diez
aperturas; $19$ ensayos para el \qty{95}{\%}; unos $4\times 10^{4}$
patrones en el \omterm{def:b3:learning-memory:hippocampus}{CA3} de la rata.
\end{solution}
